Spectroscopic binary star systems are those systems in which individual stars of the system can't be resolved visually due to their low angular separation. But we can detect that the system is binary due to shifts in their spectral line caused by Doppler shift. In double lined spectroscopic binaries, we can detect two spectra, one for each star. From the change in wavelength of the spectrum over a period of time, we can find the radial velocity curve.
Suppose an astronomer observed a hypothetical spectroscopic binary system having a circular orbit as shown in Figure 1. The plane of the orbit makes an angle $\theta$ with respect to the plane of the sky. Star 1 (mass $m_1$) is located on the outer orbit and Star 2 (mass $m_2$) is located in the inner orbit. The observed radial velocity graph is shown in Figure 2.
\ioaafig{2024_TH_T10-f1.pdf}
\ioaafig{2024_TH_T10-f2.pdf}
\begin{enumerate}
\item The binary mass function is an expression which depends upon only the observable quantities such as radial velocity and period. Find the binary mass function for both star 1 and star 2.
\item Use your expression in (1) to estimate the masses for the stars. Is your estimate exact, a lower-limit or an upper-limit?
\end{enumerate}
\textit{Supplied examination data.} $G=6.674\times10^{-11}\,\mathrm{N\,m^2\,kg^{-2}}$ and $M_\odot=1.988\times10^{30}\,\mathrm{kg}$.
